\selectlanguage{english}
\noindent\textbf{Chapter VI, manuscript pages 138--204: operations, fibres, and the choice of equivalences.}

These pages bring the construction of $\HOT_W$ to a point where it can support a calculus of operations \citep{Grothendieck1990VI}. The principal achievement is an adjunction. For a functor $f:S'\to S$, composition of structural maps already gives a direct-image functor on categories over the base. After localization, this becomes the left adjoint $f_!$ of the inverse-image functor $f^*$. The work lies in making the inverse image homotopically meaningful, identifying the unit and counit, and proving that their two composites are identities. The ensuing discussion of products, base change, and locally constant objects repeatedly returns to the choice of a concrete representative and to the changes of representative that survive the intended operation.

The notation needs a fixed orientation. We follow the manuscript's presheaf convention. A split category fibred over $S$ represents a functor $S^{\mathrm{op}}\to\Cat$. For $X\to S$, the comma replacement $\Phi_S(X)$ has fibre $s\backslash X$ at $s$; its unit $i_X:X\to\Phi_S(X)$ is a right-fine equivalence. Consequently
\[
 f^*\hot_{S,W}(X)
 \simeq\hot_{S',W}\bigl(\Phi_S(X)\times_S S'\bigr),
 \qquad
 f_!\hot_{S',W}(X')=\hot_{S,W}(X').
\]
The second formula remembers the change of structural map. It does not assert that the representative remains fibred or smooth over the new base. This is the reason the first formula cannot generally be replaced by strict pullback of $X$.

The adjunction proof on pages~138--146 is unusually revealing. The unit is the $S'$-map induced by $i_{X'}$ and the structural map $X'\to S'$. For the counit one first tries to invert the pullback of $i_X$; the attempt fails because that pullback need not belong to the relevant class of equivalences. The correct formula is
\[
 \beta_X=(i_X)^{-1}\,\pr_1^{\Phi_S(X)}.
\]
Its inverse is taken in the localization over $S$, where $i_X$ is known to be invertible. One triangle identity is then the definition of the unit, and the other is naturality of $i$. Our commentary writes out this reduction, with the domains and codomains kept visible. The later passage to locally constant objects is a formal adjunction argument between reflective subcategories, once the reflection and the appropriate commutation squares are available.

\medskip\noindent\textbf{Products and the operation attached to a representative.}

The next issue is more subtle than the construction of a multiplication law might suggest. Two smooth representatives have a strict product over $S$ that computes the product in $\HOT_W(S)$. With arbitrary representatives one must use
\[
 \hot_{S,W}(X)\times_S\hot_{S,W}(Y)
 \simeq\hot_{S,W}\bigl(\Phi_S(X)\times_S\Phi_S(Y)\bigr).
\]
The first corollary on page~152 claims that one smooth factor suffices to omit both replacements. Grothendieck retracts it immediately. Over the walking arrow $0\to1$, the smooth inclusion of $0$ and the terminal-point inclusion of $1$ have empty strict intersection, whereas replacing the latter by its comma fibration changes the calculation. This small example determines the correct scope of every later product formula.

What remains for an arbitrary $X\to S$ is the endofunctor $\operatorname{multhot}_{X/S}$, obtained by multiplying smooth representatives by this particular $X$. It gives the unconditional identity
\[
 f_!f^*\simeq\operatorname{multhot}_{S'/S}.
\]
An identification with multiplication by the localized object $\hot_{S,W}(S')$ requires a further argument. The distinction concerns functorial dependence on the first representative, rather than a failure of associativity or of the usual coherence of products. We return to it at the end of the chapter, where a proof using the Grothendieck construction extends the calculation for weak $W$-fibrations. The proof and its precise assumptions are given in the development block, independently of an appeal to a general projection formula.

The formula for a left Kan extension at an object is obtained in the same concrete way. Its value at $s$ is represented by $s\backslash X'$, rather than by the strict fibre $X'_s$. The comma square carries its canonical natural transformation; it is not in general a strictly commutative square. With that transformation and the presheaf variance made explicit, the base-change mate has the standard orientation, and the computation proves it invertible. This supplies the pointwise left-Kan-extension axiom. The smooth/proper base-change theorem then follows by choosing representatives on which strict pullback is legitimate. The initial detour through a general derivator theorem is incomplete because it invokes pointwise detection of isomorphisms. The direct argument on pages~168--169 removes that dependence.

\medskip\noindent\textbf{Why four localizations appear.}

The revision beginning on page~171 is structural. Grothendieck separates the basic localizer axioms from the additional principle that a total equivalence between weak fibrations is an equivalence on every fibre. For a map $u:X\to Y$ over $S$, the classes used in the discussion are
\[
\begin{array}{ll}
 W_S &: u\in W\text{ on the total categories},\\
 W_S^{\mathrm f} &: u_s\in W\text{ for every strict fibre},\\
 W_S^{\mathrm d} &: s\backslash u\in W\text{ for every }s,\\
 W_S^{\mathrm g} &: u/s\in W\text{ for every }s.
\end{array}
\]
On weak fibrations, the fine conditions agree; their implication to total equivalence is available under the basic axioms. The converse is W(4b). Without it, the fine localization of weak fibrations and their coarse localization cannot be identified merely because both admit useful replacements. This yields $\HOT_W$, its opposite-variance counterpart, $\HOT_W^{\mathrm{lc}}$, and $\HOT_W^{\mathrm{lc}!}$, with a localization $\lambda$ between the last two.

The path functor $C$ exposes the difficulty. Its unit is a total equivalence and its values are perfect over the base. These properties give coarse equivalences of presentations. They do not, by themselves, prove that the unit is fine on a weak fibration. W(4b) supplies exactly that step, identifies the two locally constant theories, and allows the reflection formulas to be read as genuine adjunctions with fully faithful right adjoints. The audit keeps this hypothesis attached to each use of that conclusion. A calculation with action groupoids makes the distinction visible on morphisms: coarse localization can identify maps that remain distinct in the fine theory.

This is also the place to separate historical uncertainty from present knowledge. The minimality theorem for the usual homotopy localizer, in the modern basic-localizer framework, is due to \citet{Cisinski2004Minimal}. It answers the proposed search for a smaller basic localizer; it does not identify every generalized coefficient theory with ordinary homotopy theory. The manuscript's varying axioms and its changing use of ``fibration'' must still be translated before a later theorem is applied.

\medskip\noindent\textbf{Alternating replacement and the questions it answers.}

One application of $\Phi$ cannot produce all the desired geometry. The fork-shaped base with two incomparable terminal branches gives a particularly strong obstruction: for a nondegenerate localizer, the comma replacement of a nonempty category over that base cannot be proper. The commentary verifies the categorical retraction in the preceding example, repairs the central object of the printed diagram, and gives the missing last step in the emptiness argument. A separately marked extension treats a class of partially ordered bases by the same mechanism. These arguments concern properness of a particular replacement, not the nonexistence of any satisfactory presentation.

Grothendieck then alternates the two comma constructions and takes the colimit of the resulting sequence. The odd and even subsequences respectively provide smooth and proper presentations. The proof uses W(6 bis): closure of $W$ under colimits of sequences of arrows. Its historical formulation does not assert arbitrary filtered closure. The limiting object is perfect, and its unit is fine when the original object is a weak fibration. This gives a fine equivalence between the weak-fibration and perfect-object presentations without using W(4b). It does not identify fine and coarse equivalences.

The later theorem of \citet[Corollary~4.2.22]{Cisinski2006} makes W(6 bis) automatic for basic localizers. With the identification $W_\omega=W'_\omega=W''_\omega=W_\infty$ established in the commentary, the replacement therefore answers the perfect-presentation question affirmatively under the standing axioms. The same fine units transfer the original restricted W(4b$'$) test to all weak fibrations. These conclusions concern the relative localization axiom W(2); they do not silently extend a colimit theorem to every weaker axiom system considered in the manuscript.

The practical proposition on pages~196--197 combines this construction with restricted fibre detection. Its last adjunction verification is expressly postponed. We supply it by first transporting fibre detection along the functorial replacement and then comparing the resulting unit with the established $C$-reflection. The later Chapter XVI confirms the role of replacement in relating the fibration axioms; the brief status block locates its theorem and keeps its historical hypotheses visible. No priority claim is attached to the elementary arguments supplied here.

\medskip\noindent\textbf{The reach of the construction.}

For the ordinary homotopy theory of spaces, the contemporary description of locally constant diagrams uses homotopy Kan extensions and localization of diagram categories; \citet{Cisinski2009LocallyConstant} gives a precise treatment. The model-category construction of derivators supplies a standard framework \citep{Cisinski2003,Groth2013}. More generally, Der~2 holds for every basic localizer \citep[Propositions~3.1.30 and~4.2.7]{Cisinski2006}. The commentary after manuscript p.~159 explains the comparison which permits that conclusion. The unrestricted strict-bifibration presentation over an arbitrary base remains outside what we have established: a free bifibration alone does not provide a fine equivalence from the given weak fibration. Over groupoid bases we give the complete comparison directly.

The surviving final recap deliberately lists Der~1, Der~$4'$, Der~$5'$, and the appropriate W-relative form of Der~3b. Der~2 has been withdrawn from the earlier list pending a proof; its later resolution does not supply a missing argument to the historical text. The audit also preserves the missing manuscript page~170 and the incomplete final base-change paragraph. What can be proved from the surviving material is supplied in commentary; the missing text is not reconstructed.
